% Plan: development/outline.md, section 6, item B.1.
% Sources: development/dossiers/lnts-lukvle10.md and its critique,
% development/reviews/sol-lnts-theorem3.md, development/decision-register.md
% (LL-01 to LL-15, N-01 to N-07, N-28), data/numbers.json.
% Round-1 revision (adjudication G5): one certificate box per result, the
% verification-record table merged into the boxes.
\subsection{The \inst{lnts} instances and \inst{lukvle10}}\label{app:lnts}

Both certificates are Lagrangian splits (\cref{sec:split}).
For \inst{lnts}, this subsection adds to the proof of \cref{thm:lnts-opt} (\cref{app:splitproofs-lnts}) the closed forms of the weights and the computation of the enclosures; for \inst{lukvle10}, it proves \cref{thm:lukvle10-bracket}.

% ---------------------------------------------------------------------------
\subsubsection{\inst{lnts}: stored model}\label{app:lnts-model}

The instances \inst{lnts50}, \inst{lnts100}, \inst{lnts200} and \inst{lnts400} are the GAMS Model Library model \texttt{lnts} \citep{dolan2001-benchmarking-optimization-software-with-cops,library2026-cops-and-gams-source-models}, stated in \cref{sec:split-affine}.
Each has $5N+6$ continuous variables and $4N$ rows, the dynamics written componentwise; each velocity row is the product of $100f(\theta_i)+100f(\theta_{i+1})$, the number $-.5$ and $h$, with $f\in\{\cos,\sin\}$.
Several fixed values are written as an upper bound $0$ together with the OSIL default lower bound $0$; a reader that took a missing lower bound as $-\infty$ would read a different model.
The rows involve only the data $100$, $-.5$, $45$, $5$ and $0$; the angle bound $\bar\theta=1.5707963267949$ is the only datum that is neither an integer nor a dyadic rational.

% ---------------------------------------------------------------------------
\subsubsection{\inst{lnts}: weights, enclosure and stored points}\label{app:lnts-proof}

We use the notation of \cref{sec:split-affine}: the weights $w_j$ and $c_j$, $r_j=c_j/w_j-N/2$, $\alpha(h)=9/(20h)$, $\beta(h)=1/(20h^2)$, and $C$, $D$ and $g$ as in \eqref{eq:lnts-g}.
The proofs of \cref{lem:lnts-elim,thm:lnts-opt} use the following closed forms and reflection identity.

\begin{lemma}[closed forms of the weights]\label{lem:lnts-weights}
Let $\theta\in\R^{N+1}$ and $h\in\R$, and let the states follow from the dynamics rows and the zero initial states.
\begin{enumerate}
\item[(a)] For $0\le k\le N$, $v^x_k=100h\sum_jW_{kj}\cos\theta_j$ and $v^y_k=100h\sum_jW_{kj}\sin\theta_j$, where $W_{kj}=\tfrac12\bigl([\,j\le k-1\,]+[\,1\le j\le k\,]\bigr)$; in particular $W_{Nj}=w_j$.
Further, $p^y_N=100h^2\sum_jc_j\sin\theta_j$ with $c_j=\sum_kw_kW_{kj}$, and these sums equal the closed forms $c_0=(2N-1)/4$, $c_j=N-j$ for $0<j<N$, and $c_N=\tfrac14$.
\item[(b)] $c_{N-j}=Nw_j-c_j$ and $r_{N-j}=-r_j$ for all $j$.
Hence $\sum_jw_jr_j=0$ and $r_{N/2}=0$, and $r_j\ne0$ for $j\ne N/2$.
\end{enumerate}
\end{lemma}

\begin{proof}
(a) Summing the vertical velocity rows gives $v^y_k=50h\sum_{i<k}(\sin\theta_i+\sin\theta_{i+1})$.
The index $j=i$ runs over $0\le j\le k-1$ and the index $j=i+1$ over $1\le j\le k$, which gives $W_{kj}$; the horizontal velocity is the same with cosines.
For $k=N$, $W_{N0}=W_{NN}=\tfrac12$ and $W_{Nj}=1$ otherwise.
Summing the vertical position rows and using $p^y_0=0$ gives $p^y_N=\tfrac h2\sum_{i<N}(v^y_i+v^y_{i+1})=h\sum_kw_kv^y_k$, hence the stated form with $c_j=\sum_kw_kW_{kj}$.
For $j=0$, $W_{k0}=\tfrac12$ for $k\ge1$ and $W_{00}=0$, so $c_0=\tfrac12\sum_{k\ge1}w_k=\tfrac12(N-\tfrac12)=(2N-1)/4$.
For $0<j<N$, $W_{kj}=\tfrac12([\,k\ge j+1\,]+[\,k\ge j\,])$, so $c_j=\tfrac12\bigl((N-j-\tfrac12)+(N-j+\tfrac12)\bigr)=N-j$.
For $j=N$, $W_{kN}=\tfrac12[\,k=N\,]$, so $c_N=w_N/2=\tfrac14$.

(b) Directly, $c_N=\tfrac14=N/2-c_0$, $c_{N-j}=j=N-c_j$ for $0<j<N$, and $c_0=N/2-c_N$.
Since $w_{N-j}=w_j$, it follows that $r_{N-j}=(Nw_j-c_j)/w_j-N/2=-r_j$.
The terms of $\sum_jw_jr_j$ cancel in pairs, and $N$ is even, so $r_{N/2}=-r_{N/2}=0$.
Finally $r_0=-r_N=(N-1)/2$ and $r_j=N/2-j$ for $0<j<N$, which vanishes only at $j=N/2$.
\end{proof}

The bound $\bar\theta$ enters the proof of \cref{thm:lnts-opt} only through the admissibility of $\theta^*$, which holds with room to spare: the certified brackets of $\nu^*$ give $|t_j|=\nu^*|r_j|<1.41$, so $|\theta^*_j|<0.96$.
The multiplier $\mu^*=-\nu^*N/2$ reflects the symmetry $\theta_j\mapsto-\theta_{N-j}$, which by \cref{lem:lnts-weights}(b) preserves the three terminal equations \eqref{eq:lnts-terminal} at every fixed $h$ and so maps the exactly feasible points onto themselves.

\paragraph{Computation of the enclosure.}
Both codes compute the rationals $a<b$ and the bounds on $C(a)$ and $C(b)$ of the proof of \cref{thm:lnts-opt} with Python integers and fractions only.
The first code uses the closed forms of \cref{lem:lnts-weights} and encloses $\sqrt q$ for a rational $q=p/r$ by $[m,m+1]/(r\,2^{600})$ with the integer square root $m=\lfloor\sqrt{pr\,2^{1200}}\rfloor$.
The second code derives the weights by symbolic forward recursion, encloses each $(1+t^2)^{-1/2}$, where $1+t^2=p/d$, by $k/M$ and $(k+1)/M$ with $M=10^{120}$ and $k^2p\le dM^2<(k+1)^2p$, and bisects $300$ times from $[0,4/N]$, accepting a half only when the whole enclosure of $g$ there has a strict sign.
Its rational enclosures of $\vstar$ give the displays; \cref{tab:lnts-enclosure} rounds them outward to 40 decimals.

\begin{table}[ht]
\centering
\small
\caption{The optimal values of \cref{thm:lnts-opt}.
The root $\nu^*$ is shown truncated, for orientation.
The last column is the 40-decimal lower end of the enclosure of $\vstar=Nh^*$; the upper end is larger by $10^{-40}$.
The underlying rational enclosures are narrower than \sci{1.11}{-91}.}
\label{tab:lnts-enclosure}
\begin{tabular}{@{}lll@{}}
\toprule
instance & $\nu^*$ & lower end of the enclosure of $\vstar$\\
\midrule
\inst{lnts50}  & $0.05640028332758\ldots$ & $0.5546687649386788986220922339726517468015$\\
\inst{lnts100} & $0.02817764593774\ldots$ & $0.5545954011669111610017828739043968162226$\\
\inst{lnts200} & $0.01408601756576\ldots$ & $0.5545770161030836720556252926039795506005$\\
\inst{lnts400} & $0.00704265847933\ldots$ & $0.5545724137006871088173911331507631508456$\\
\bottomrule
\end{tabular}
\end{table}

\begin{certbox}[Certificate for \cref{thm:lnts-opt}]
\certitem{Inputs}{OSIL files (SHA-256 \texttt{3d8234b3}, \texttt{284bddc3}, \texttt{c2dde9ef}, \texttt{96eb4456}).}
\certitem{Computation}{model structure and weights checked; 300 bisections of $g$ on strict signs of enclosures; integer square roots; \arith{E}.}
\certitem{Data reading}{exact, by two separately written readers; $\bar\theta$ enters only through $\pi/2<\bar\theta$.}
\certitem{Implementations}{two separately written integer-only codes, each with its own OSIL reader; the second gives the displays (brackets of width below $10^{-90}$), the first encloses the optima to width below $5\cdot10^{-61}$; a third code, in 130-digit \texttt{mpmath} intervals, is consistent.}
\certitem{Evidence level}{\evid{hand}+\evid{stored}.}
\certitem{Replay}{Tier~1; seconds per instance.}
\end{certbox}

\begin{remark}[the stored points are strictly suboptimal]\label{rem:lnts-stored}
\Cref{app:points} proves by a Krawczyk test on the three terminal equations (\cref{thm:pts-krawczyk}) that a second set of exactly feasible points of the \inst{lnts} models exists.
These points fix $\theta_1,\dots,\theta_{N-1}$ at 40-digit decimals, and their middle control $\theta_{N/2}$ is a nonzero rational of magnitude below $10^{-110}$.
They are strictly suboptimal.
Indeed, let $x$ be exactly feasible with controls $\theta$ and step $h$, and put $t_j=\nu^*r_j$.
By \eqref{eq:lnts-terminal} and $\mu^*w_j+\nu^*c_j=w_jt_j$,
\[
\alpha(h)+\nu^*\beta(h)=\sum_jw_j\bigl(\cos\theta_j+t_j\sin\theta_j\bigr)=\alpha(h^*)+\nu^*\beta(h^*)-\sum_j\delta_j,
\]
where $\delta_j=w_j\bigl(\sqrt{1+t_j^2}-\cos\theta_j-t_j\sin\theta_j\bigr)\ge0$ by the Cauchy--Schwarz inequality, and the second equality uses $S(\mu^*,\nu^*)=\sum_jw_j\sqrt{1+t_j^2}=\alpha(h^*)+\nu^*\beta(h^*)$ from the proof of \cref{thm:lnts-opt}.
Since $t_{N/2}=0$ and $w_{N/2}=1$, $\delta_{N/2}=1-\cos\theta_{N/2}$, which is positive for these points.
So $\alpha(h)+\nu^*\beta(h)<\alpha(h^*)+\nu^*\beta(h^*)$, and $h>h^*$ because the left side is strictly decreasing in $h$.
The objective enclosures of these points from the Krawczyk test (110-digit interval arithmetic) are narrower than $3\cdot10^{-109}$.
Their stored 60-digit enclosures of $h$, multiplied by $N$, overlap the rational enclosure of $\vstar$ and show that their objective values exceed $\vstar$ by less than $2\cdot10^{-58}$.
This closeness is a consistency check between two separate existence proofs; it does not make these points optimal.
\end{remark}

\paragraph{A listed point below the optimum.}
MINLPLib's point p1 of \inst{lnts50} has the step $h=0.011093375297913$, so its objective $Nh=0.55466876489565$ lies at least \sci{4.30}{-11} below $\vstar$, and its maximal row violation is \sci{9.1}{-10}.
By \cref{lem:sem-enclosure}(b) it is not exactly feasible; it is feasible only within MINLPLib's tolerance (category~A, \cref{tab:claims}).
The points p1 of \inst{lnts100}, \inst{lnts200} and \inst{lnts400} evaluate above $\vstar$.
The mechanism of \cref{thm:lnts-opt} is a classical sufficiency argument; to our knowledge, within the search of \cref{app:literature}, the attaining discrete control, the scalar equation $g(\nu)=0$ and the rigorous enclosure are new (the local values of COPS~2.0 agree with our optima to six digits).

% ---------------------------------------------------------------------------
\subsubsection{\inst{lukvle10}: stored model}\label{app:lukvle10-model}

The instance \inst{lukvle10} is stated in \cref{sec:split-window}; it is an artificial test problem that originates from problem~5.10 of \citet{luksan1999-test-problems-for-unconstrained-optimization}.
The OSIL file writes the objective terms as \texttt{power}(\texttt{square}($x_a$), \texttt{square}($x_b$)$+1$) and the rows as $-x_j+3x_{j+1}-2x_{j+2}-2x_{j+1}^2=-1$.
By \cref{def:sem-model}(iii), $(a^2)^{e}=\exp(e\ln a^2)$ for $a\neq0$ and $0^{e}=0$, because the exponent $e=b^2+1$ is positive.
Thus $F$ is defined and nonnegative on $\R^2$.
It is also continuously differentiable: near $a=0$ the term $(a^2)^{b^2+1}=|a|^{2(b^2+1)}$ and its partial derivatives tend to $0$, because the exponent $2(b^2+1)$ is at least $2$, and the same holds for the other term near $b=0$.
All data of the model are integers, so the three data readings of \cref{sec:semantics-model} give the same model.

Row $j$ is linear in $x_{j+2}$ with coefficient $-2$ and defines $x_{j+2}=\Phi(x_j,x_{j+1})$ with $\Phi(u,v)=(1-u+3v-2v^2)/2$, so the exactly feasible points correspond one to one to the seeds $(x_0,x_1)\in\R^2$.
The map $(u,v)\mapsto(v,\Phi(u,v))$ has the saddle fixed point $(p,p)$ with $p=-1/\sqrt2$ and eigenvalues about $2.731$ and $0.183$ (\cref{rem:lukvle10-seeds}).
The best known points follow this saddle for most indices, where each pair costs $F(p,p)=2^{-1/2}$, and have short transients at both ends.

The CUTEst SIF file records the value \texttt{SOLTN 3.52237E+02}, which lies at least \sci{1.02}{-3} below our certified dual bound.
No exactly feasible point attains it (\cref{lem:sem-enclosure}(b)); its origin is not documented (\cref{tab:claims}).

% ---------------------------------------------------------------------------
\subsubsection{\inst{lukvle10}: partial Lagrangian, window and proof}\label{app:lukvle10-proof}

We dualize the first $m=994$ rows and keep rows $994$--$997$.
For $\lambda\in\R^m$ put $\lambda_j=0$ for $j\notin[0,m)$ and, for $k=0,\dots,995$,
\begin{equation}\label{eq:lukvle10-coef}
\beta_k=-\lambda_k+3\lambda_{k-1}-2\lambda_{k-2},\qquad q_k=-2\lambda_{k-1}.
\end{equation}
In particular $q_0=0$, $\beta_{994}=3\lambda_{993}-2\lambda_{992}$, $q_{994}=-2\lambda_{993}$, $\beta_{995}=-2\lambda_{993}$ and $q_{995}=0$.
Define
\begin{align*}
\ell_i(a,b)&=F(a,b)+\beta_{2i}a+q_{2i}a^2+\beta_{2i+1}b+q_{2i+1}b^2, \qquad i=0,\dots,496,\\
\tau(a,b)&=F(a,b)+F\bigl(P(a,b)\bigr)+F\bigl(P^2(a,b)\bigr)+\beta_{994}a+q_{994}a^2+\beta_{995}b,
\end{align*}
where $P(a,b)=\bigl(u,\Phi(b,u)\bigr)$ with $u=\Phi(a,b)$ maps a pair $(x_{2i},x_{2i+1})$ of an exactly feasible point to the next pair.

\begin{lemma}[partial Lagrangian]\label{lem:lukvle10-partial}
For every $\lambda\in\R^{994}$, every exactly feasible point $x$ of \inst{lukvle10} satisfies
\[
f(x)\ \ge\ \sum_{j<994}\lambda_j+\sum_{i=0}^{496}\inf_{\R^2}\ell_i+\inf_{\R^2}\tau .
\]
\end{lemma}

\begin{proof}
For an exactly feasible $x$, $c_j(x)=0$ for all $j$, where $c_j(x)=1-x_j+3x_{j+1}-2x_{j+2}-2x_{j+1}^2$, so $f(x)=f(x)+\sum_{j<m}\lambda_jc_j(x)$.
Row $j$ contributes $\lambda_j-\lambda_jx_j+3\lambda_jx_{j+1}-2\lambda_jx_{j+2}-2\lambda_jx_{j+1}^2$.
Collecting the coefficient of each $x_k$ and of each $x_k^2$ gives $\beta_k$ and $q_k$ in \eqref{eq:lukvle10-coef}; the dualized rows reach only up to $x_{995}$.
The kept rows $994$--$997$ give $(x_{996},x_{997})=P(x_{994},x_{995})$ and $(x_{998},x_{999})=P^2(x_{994},x_{995})$.
Hence
\[
f(x)=\sum_{j<m}\lambda_j+\sum_{i=0}^{496}\ell_i(x_{2i},x_{2i+1})+\tau(x_{994},x_{995}),
\]
and each term is at least its infimum over $\R^2$.
\end{proof}

This is \cref{lem:split-bound}(d) with a window of three pairs, the row-form analogue of \cref{prop:split-window}; the window is two-dimensional because its entry pair $(x_{994},x_{995})$ determines its interior through $P$.

\begin{lemma}[reduction to a box]\label{lem:lukvle10-coercive}
Let $\varphi(t)=t^2$ for $|t|\ge1$ and $\varphi(t)=0$ otherwise, and for $q>-1$ and $\beta\in\R$ let $\psi(t)=\varphi(t)+qt^2+\beta t$.
\begin{enumerate}
\item[(a)] $F(a,b)\ge\varphi(a)+\varphi(b)$; hence $\ell_i(a,b)\ge\psi_{2i}(a)+\psi_{2i+1}(b)$ and $\tau(a,b)\ge\psi_{994}(a)+\psi_{995}(b)$, where $\psi_k$ has the coefficients $(q_k,\beta_k)$.
\item[(b)] $\inf\psi\ge m_\psi:=\min\bigl(-|q|-|\beta|,\ -\beta^2/(4(1+q))\bigr)$.
\item[(c)] Let $\ell$ be one of the functions $\ell_i$ or $\tau$, with coefficient pairs $(q_a,\beta_a)$ and $(q_b,\beta_b)$, and let $U\ge\inf_{\R^2}\ell$.
If $R_a\ge\max\bigl(1,|\beta_a|/(2(1+q_a))\bigr)$ and $(1+q_a)R_a^2-|\beta_a|R_a+m_{\psi_b}>U$, and symmetrically for $R_b$, then
$\inf_{\R^2}\ell=\min\{\ell(a,b):|a|\le R_a,\ |b|\le R_b\}$, and the minimum is attained.
\end{enumerate}
\end{lemma}

\begin{proof}
(a) Both terms of $F$ are nonnegative, and $(a^2)^{b^2+1}\ge a^2$ when $a^2\ge1$, since the exponent is at least $1$; the same holds for the second term.
For $\tau$, drop the nonnegative terms $F\circ P$ and $F\circ P^2$.
(b) For $|t|<1$, $\psi(t)=qt^2+\beta t\ge-|q|-|\beta|$.
For $|t|\ge1$, $\psi(t)\ge(1+q)t^2-|\beta||t|\ge-\beta^2/(4(1+q))$.
(c) The function $s\mapsto(1+q_a)s^2-|\beta_a|s$ is nondecreasing for $s\ge|\beta_a|/(2(1+q_a))$.
So for $|a|\ge R_a\ge1$ and every $b$, (a) and (b) give $\ell(a,b)\ge(1+q_a)R_a^2-|\beta_a|R_a+m_{\psi_b}>U\ge\inf\ell$; the same holds for $|b|\ge R_b$.
Hence the infimum over $\R^2$ equals the infimum over the box, and $\ell$ is continuous on the compact box.
\end{proof}

The lemma needs only $q_k>-1$, that is, $\lambda_{k-1}<\tfrac12$; the multipliers below give $q_k\in[-0.8156,0]$ for all $k$.
Part~(b) bounds an infimum because $\psi$ is discontinuous at $|t|=1$.

\paragraph{Multipliers.}
Any $\lambda$ gives a valid bound, provided the same vector enters the coefficients $\beta_k,q_k$ and the sum $\sum_j\lambda_j$.
The certificate uses least-squares multipliers of the KKT system at MINLPLib's point p5, computed in binary64 (residual \sci{4.6}{-14}).
It then replaces $\lambda_{30},\dots,\lambda_{960}$, whose spread is \sci{2.4}{-14}, by the single binary64 value $\bar\lambda=\lambda_{495}=0.4077978179563417$.
Pair problem $i$ depends on $\lambda_{2i-2},\dots,\lambda_{2i+1}$, so this makes the 464 pair problems $i=16,\dots,479$ identical, with $\beta=0$ and $q=-2\bar\lambda$, and one search bounds all of them.
The value $\bar\lambda$ is close to $(3-\ln2)/(4\sqrt2)$, the multiplier for which $F(a,b)-2\bar\lambda(a^2+b^2)$ is stationary at the saddle pair $(-1/\sqrt2,-1/\sqrt2)$.
All multipliers lie in $[0.2341,0.4078]$, and they enter the computation as exact binary64 numbers.

\paragraph{Branch and bound.}
Each of the 34 distinct pair problems and the window problem $\tau$ is bounded below by a best-first interval branch and bound on the box of \cref{lem:lukvle10-coercive}(c).
The radii are chosen with $U$ equal to the interval upper end of the function at the KKT pair plus $1$, and the inequality of \cref{lem:lukvle10-coercive}(c) is checked in interval arithmetic; the margin $1$ also covers the rounding of $|\beta|$ to binary64 in the choice of the radii.
All arithmetic is \texttt{mpmath} interval arithmetic with a 64-bit mantissa (\arith{I}).
The range of $(a^2)^{e}$ over a box follows from monotonicity (increasing in the base; in $e$, increasing for bases at least $1$ and decreasing otherwise), with endpoint values $\exp(e\ln a^2)$, and $0$ at base $0$.
The lower bound on a box $X$ with centre $c$ is the larger of the natural interval extension and the mean-value form $\ell(c)+\nabla\ell(X)\cdot(X-c)$, which is valid because $\ell$ is continuously differentiable and $c\in X$.
The interval gradient comes from forward-mode differentiation, through $P$ and $P^2$ for $\tau$; it is used only when no argument interval of $F$ contains $0$, so that the logarithms are defined.
The incumbent is the interval upper end of $\ell$ at some point (initially the KKT pair, later box centres), so a box whose lower bound exceeds it contains no minimizer and is discarded.
The search stops when the incumbent exceeds the smallest lower bound of the open boxes by at most $10^{-12}$ (pairs) or $10^{-9}$ (window) and returns that smallest lower bound.
The discarded and open boxes partition the root box, and a box that contains a minimizer is never discarded, so the minimum is at least the returned value.
\Cref{tab:lukvle10-bnb} summarizes the searches.

\begin{table}[ht]
\centering
\small
\caption{Interval branch and bound for \inst{lukvle10}.
Half-widths are those of the boxes of \cref{lem:lukvle10-coercive}(c); UB$-$LB is the final incumbent minus the returned lower bound, and the last column counts each pair problem with its multiplicity.
In total the searches evaluated 329{,}233 boxes; the trees are not stored.}
\label{tab:lukvle10-bnb}
\begin{tabular}{@{}llllll@{}}
\toprule
subproblems & searches & half-widths & boxes per search & max.\ UB$-$LB & sum of UB$-$LB\\
\midrule
pairs $0$--$15$ & 16 & $1.44$--$3.043$ & 2{,}255--9{,}365 & $1.0\cdot10^{-12}$ & $1.53\cdot10^{-11}$\\
pairs $16$--$479$ & 1 (for 464) & $3.043$ & 9{,}471 & $8.4\cdot10^{-13}$ & $3.89\cdot10^{-10}$\\
pairs $480$--$496$ & 17 & $2.998$--$3.043$ & 4{,}473--9{,}415 & $1.0\cdot10^{-12}$ & $1.59\cdot10^{-11}$\\
window $\tau$ & 1 & $4.99$, $2.28$ & 92{,}551 & $1.0\cdot10^{-9}$ & $1.0\cdot10^{-9}$\\
\bottomrule
\end{tabular}
\end{table}

\begin{proof}[Proof of \cref{thm:lukvle10-bracket}]
\emph{Lower bound.}
Let $x$ be exactly feasible.
By \cref{lem:lukvle10-partial}, $f(x)\ge\sum_j\lambda_j+\sum_i\inf\ell_i+\inf\tau$ for the multipliers above.
By \cref{lem:lukvle10-coercive}, each infimum is a minimum over the box of \cref{tab:lukvle10-bnb}, and by the branch-and-bound argument it is at least the returned lower bound.
The returned bounds are logged to 20 significant digits, and each logged value is reduced by $10^{-18}$, so that it stays below the returned bound.
Summing $\sum_j\lambda_j$ and these values (the middle one with multiplicity 464) in interval arithmetic gives the lower end $352.2380254050784556\ldots$; an exact rational recomputation from the logged values gives $352.2380254050784557\ldots$.
Both exceed the display $352.2380254050784$ by more than $5\cdot10^{-14}$, so $\vstar\ge352.2380254050784$.
This step assumes that \texttt{mpmath}~1.3.0 interval arithmetic encloses the arithmetic operations, integer powers, $\exp$ and $\ln$ correctly (trust term \trust{mp}, \cref{app:repro-setup}).

\emph{Upper bound.}
The point $x^\star$ has seeds $x_0,x_1$ equal to KKT values rounded to 640 decimals, and $x_2,\dots,x_{999}$ are defined by $x_{j+2}=\Phi(x_j,x_{j+1})$.
It is rational, it satisfies all 998 rows exactly, and every variable is free.
Its coordinates are enclosed in intervals of radius below \sci{5}{-61}, and $\min_k|x_k|>0.26$.
Its objective lies in $[y,\,y+10^{-40}]$ with
\[
y=352.2380254064956226308712710293664647979978,
\]
so $\vstar\le f(x^\star)\le352.2380254064957$.
Hence $\gapabs\le\sci{1.42}{-9}$ and $\gaprel\le\sci{4.03}{-12}$.
\end{proof}

The point $x^\star$ is defined by its seeds and the rows; its denominators roughly square at each step, so its coordinates cannot usefully be written out (\cref{rem:lukvle10-seeds}).
Two implementations that use only integer arithmetic, with explicit series remainder bounds for $\exp$ and $\ln$, reproduce the enclosure (\cref{app:points}), so the primal claim does not depend on \texttt{mpmath}.

\begin{remark}[what the gap means]\label{rem:lukvle10-gap}
The certified gap, at most \sci{1.42}{-9}, is dominated by the summed branch-and-bound tolerances (at most \sci{4.20}{-10} for the 497 pairs and \sci{9.98}{-10} for the window); they agree to about $3.5\cdot10^{-16}$.
Global optimality of the exactly feasible point is not proved.
The incumbents of the searches never fell below the values of the pair Lagrangians at the KKT pairs, and floating-point multistart searches found no lower point (numerical evidence).
In floating-point experiments, dualizing all 998 rows lost about $0.09$ in the last pairs, where the trajectory leaves the saddle and the pair Lagrangians are nonconvex; keeping rows $994$--$997$ removes this loss up to the tolerances above.
A longer window would stay two-dimensional, but its interval enclosures would widen by about $2.73$ per step.
\end{remark}

Partial Lagrangian relaxation and interval branch and bound with natural and mean-value enclosures are classical \citep{neumaier2004-complete-search-in-continuous-global}; the contribution is the certificate for this model.

\begin{certbox}[Certificate for \cref{thm:lukvle10-bracket}]
\certitem{Inputs}{OSIL file (SHA-256 \texttt{1eb1236d}); binary64 multipliers (\texttt{df968aac}); seeds (\texttt{767d1e25}) and coordinate box (\texttt{7a106e83}) of the primal point.}
\certitem{Computation}{dual: 35 interval branch-and-bound searches and their sum, \arith{I} (\texttt{mpmath}~1.3.0, 64-bit; $+,-,\times,\div$, integer powers, $\exp$, $\ln$), the sum also exactly from the logged bounds; primal: objective enclosed in integer arithmetic, \arith{E}.}
\certitem{Data reading}{integer data, so every reading gives the same model; \texttt{power} trees matched as strings.}
\certitem{Implementations}{the second code certifies the displayed dual; the first code (own reader, multipliers, radii and search, 30 digits) certifies the weaker $352.238025369202$; both trust \texttt{mpmath}'s interval $\exp$ and $\ln$. Two integer-only codes certify the primal enclosure.}
\certitem{Evidence level}{dual \evid{rerun} (search trees not stored); primal \evid{stored}.}
\certitem{Replay}{Tier~1; about 300\,s for the dual searches.}
\end{certbox}
